\begin{filecontents*}{refs.bib}
@article{lichti2010propagation,
  author  = {Lichti, D. D. and O'Keefe, K. and Jamtsho, S.},
  title   = {Propagation of an Unmodeled Additive Constant in Range Sensor Observations},
  journal = {Journal of Surveying Engineering},
  year    = {2010}, volume = {136}, number = {3}, pages = {111--119},
  doi     = {10.1061/(ASCE)SU.1943-5428.0000024}
}
@inproceedings{kuang2013stratified,
  author    = {Kuang, Yubin and {\AA}str{\"o}m, Kalle},
  title     = {Stratified Sensor Network Self-Calibration from {TDOA} Measurements},
  booktitle = {21st European Signal Processing Conference (EUSIPCO 2013)},
  year      = {2013}, pages = {1--5}, publisher = {IEEE},
  note      = {verify page range}
}
@inproceedings{kuang2013complete,
  author    = {Kuang, Yubin and Burgess, Simon and Torstensson, Anna and {\AA}str{\"o}m, Kalle},
  title     = {A Complete Characterization and Solution to the Microphone Position Self-Calibration Problem},
  booktitle = {2013 IEEE International Conference on Acoustics, Speech and Signal Processing (ICASSP)},
  year      = {2013}, pages = {3875--3879}, publisher = {IEEE},
  doi       = {10.1109/ICASSP.2013.6638382}
}
@inproceedings{batstone2019robust,
  author    = {Batstone, Kenneth and Flood, Gabrielle and Beleyur, Thejasvi and Larsson, Viktor and Goerlitz, Holger R. and Oskarsson, Magnus and {\AA}str{\"o}m, Kalle},
  title     = {Robust Self-Calibration of Constant Offset Time-Difference-of-Arrival},
  booktitle = {2019 IEEE International Conference on Acoustics, Speech and Signal Processing (ICASSP)},
  year      = {2019}, pages = {4410--4414}, publisher = {IEEE},
  doi       = {10.1109/ICASSP.2019.8683221}
}
@inproceedings{crocco2012closed,
  author    = {Crocco, Marco and Del Bue, Alessio and Bustreo, Marco and Murino, Vittorio},
  title     = {A Closed Form Solution to the Microphone Position Self-Calibration Problem},
  booktitle = {2012 IEEE International Conference on Acoustics, Speech and Signal Processing (ICASSP)},
  year      = {2012}, pages = {2597--2600}, publisher = {IEEE},
  doi       = {10.1109/ICASSP.2012.6288448}
}
@article{simayijiang2015toa,
  author  = {Simayijiang, Zhayida and Burgess, Simon and Kuang, Yubin and {\AA}str{\"o}m, Kalle},
  title   = {{TOA}-Based Self-Calibration of Dual-Microphone Array},
  journal = {IEEE Journal of Selected Topics in Signal Processing},
  year    = {2015}, volume = {9}, number = {5}, pages = {791--801},
  doi     = {10.1109/JSTSP.2015.2417117}
}
@article{shah2021antenna,
  author  = {Shah, Syed Aamir Hussain and others},
  title   = {Antenna Delay Calibration of {UWB} Nodes},
  journal = {IEEE Access},
  year    = {2021}, volume = {9}, pages = {63294--63305},
  doi     = {10.1109/ACCESS.2021.3074872},
  note    = {verify full author list, volume and pages}
}
@inproceedings{shalaby2023calibration,
  author    = {Shalaby, Mohammed Ayman and Cossette, Charles Champagne and Forbes, James Richard and Le Ny, Jerome},
  title     = {Calibration and Uncertainty Characterization for Ultra-Wideband Two-Way-Ranging Measurements},
  booktitle = {2023 IEEE International Conference on Robotics and Automation (ICRA)},
  year      = {2023}, pages = {4128--4134}, publisher = {IEEE},
  doi       = {10.1109/ICRA48891.2023.10160769}
}
@article{piavanini2022selfcalibrating,
  author  = {Piavanini, Marco and Barbieri, Luca and Brambilla, Mattia and Cerutti, Mattia and Ercoli, Simone and Agili, Andrea and Nicoli, Monica},
  title   = {A Self-Calibrating Localization Solution for Sport Applications with {UWB} Technology},
  journal = {Sensors},
  year    = {2022}, volume = {22}, number = {23}, pages = {9363},
  doi     = {10.3390/s22239363}
}
@inproceedings{piavanini2022calibration,
  author    = {Piavanini, Marco and Barbieri, Luca and Brambilla, Mattia and Cerutti, Mattia and Ercoli, Simone and Agili, Andrea and Nicoli, Monica},
  title     = {A Calibration Method for Antenna Delay Estimation and Anchor Self-Localization in {UWB} Systems},
  booktitle = {2022 IEEE International Workshop on Metrology for Industry 4.0 \& IoT (MetroInd4.0\&IoT)},
  year      = {2022}, pages = {173--177}, address = {Trento, Italy}, publisher = {IEEE},
  doi       = {10.1109/MetroInd4.0IoT54413.2022.9831579}
}
@article{hamer2018selfcalibrating,
  author  = {Hamer, Michael and D'Andrea, Raffaello},
  title   = {Self-Calibrating Ultra-Wideband Network Supporting Multi-Robot Localization},
  journal = {IEEE Access},
  year    = {2018}, volume = {6}, pages = {22292--22304},
  doi     = {10.1109/ACCESS.2018.2829020}
}
@misc{nguyen2025coordinate,
  author        = {Nguyen, Tien-Dat and Nguyen, Thien-Minh and Nguyen, Vinh-Hao},
  title         = {Coordinate-Consistent Localization via Continuous-Time Calibration and Fusion of {UWB} and {SLAM} Observations},
  year          = {2025}, eprint = {2510.05992}, archivePrefix = {arXiv}, primaryClass = {cs.RO}
}
@inproceedings{grosswindhager2019snaploc,
  author    = {Gro{\ss}windhager, Bernhard and Stocker, Michael and Rath, Michael and Boano, Carlo Alberto and R{\"o}mer, Kay},
  title     = {{SnapLoc}: An Ultra-Fast {UWB}-Based Indoor Localization System for an Unlimited Number of Tags},
  booktitle = {2019 18th ACM/IEEE International Conference on Information Processing in Sensor Networks (IPSN)},
  year      = {2019}, pages = {61--72}, publisher = {IEEE},
  doi       = {10.1145/3302506.3310389}
}
@article{sidorenko2020error,
  author  = {Sidorenko, Juri and Schatz, Volker and Scherer-Negenborn, Norbert and Arens, Michael and Hugentobler, Urs},
  title   = {Error Corrections for Ultrawideband Ranging},
  journal = {IEEE Transactions on Instrumentation and Measurement},
  year    = {2020}, volume = {69}, number = {11}, pages = {9037--9047},
  doi     = {10.1109/TIM.2020.2996706}
}
@inproceedings{dotlic2018ranging,
  author    = {Dotlic, Igor and Connell, Andrew and McLaughlin, Michael},
  title     = {Ranging Methods Utilizing Carrier Frequency Offset Estimation},
  booktitle = {2018 15th Workshop on Positioning, Navigation and Communications (WPNC)},
  year      = {2018}, pages = {1--6}, publisher = {IEEE},
  doi       = {10.1109/WPNC.2018.8555763}
}
@inproceedings{neirynck2016alternative,
  author    = {Neirynck, Dries and Luk, Eric and McLaughlin, Michael},
  title     = {An Alternative Double-Sided Two-Way Ranging Method},
  booktitle = {2016 13th Workshop on Positioning, Navigation and Communications (WPNC)},
  year      = {2016}, pages = {1--4}, publisher = {IEEE},
  doi       = {10.1109/WPNC.2016.7822844}
}
@article{shalaby2024reducing,
  author        = {Shalaby, Mohammed Ayman and Cossette, Charles Champagne and Forbes, James Richard and Le Ny, Jerome},
  title         = {Reducing Two-Way Ranging Variance by Signal-Timing Optimization},
  journal       = {IEEE Transactions on Aerospace and Electronic Systems},
  year          = {2024}, eprint = {2211.00538}, archivePrefix = {arXiv},
  note          = {verify volume/number/pages on IEEE Xplore}
}
@article{albaidhani2021anchor,
  author  = {Albaidhani, Abbas and Alsudani, Ahmed},
  title   = {Anchor Selection by Geometric Dilution of Precision for an Indoor Positioning System Using Ultra-Wide Band Technology},
  journal = {IET Wireless Sensor Systems},
  year    = {2021}, volume = {11}, number = {1}, pages = {22--31},
  doi     = {10.1049/wss2.12006}
}
@article{albaidhani2019anchor,
  author  = {Albaidhani, Abbas and Morell, Antoni and Vicario, Jos{\'e} L{\'o}pez},
  title   = {Anchor Selection for {UWB} Indoor Positioning},
  journal = {Transactions on Emerging Telecommunications Technologies},
  year    = {2019}, volume = {30}, number = {6}, pages = {e3598},
  doi     = {10.1002/ett.3598}
}
@inproceedings{lutz2019visual,
  author    = {Lutz, Philipp and Schuster, Martin J. and Steidle, Florian},
  title     = {Visual-Inertial {SLAM} Aided Estimation of Anchor Poses and Sensor Error Model Parameters of {UWB} Radio Modules},
  booktitle = {2019 19th International Conference on Advanced Robotics (ICAR)},
  year      = {2019}, pages = {739--746}, publisher = {IEEE},
  doi       = {10.1109/ICAR46387.2019.8981544},
  note      = {verify page range}
}
\end{filecontents*}

\documentclass[11pt,a4paper]{article}

\usepackage[T1]{fontenc}
\usepackage{lmodern}
\usepackage[margin=2.5cm]{geometry}
\usepackage[protrusion=true,expansion=false]{microtype}
\usepackage{booktabs}
\usepackage{array}
\usepackage{enumitem}
\usepackage{csquotes}
\usepackage{xcolor}
\usepackage[hidelinks]{hyperref}

\newcommand{\verdict}[1]{\textbf{\textcolor{black!80}{#1}}}

\setlength{\emergencystretch}{2em}

\title{Novelty Check --- Frame-Dependent Antenna Delay \&\\ Inter-Anchor UWB Self-Calibration}
\author{Prior-art and novelty assessment --- prepared for\\ Prof.\ Bj\"orn Eskofier (LMU M\"unchen)}
\date{\today}

\begin{document}

\maketitle

\section*{Executive Summary}

This document is a prior-art and novelty assessment of the AutoPos UWB inter-anchor self-calibration work, prepared ahead of submission to \emph{Measurement} / \emph{IEEE TIM}. Five claimed contributions were each checked against the literature, with an active search for disconfirming prior art. The conclusion: the work has exactly one defensible novelty --- Contribution 1 (antenna delay is a frame-dependent, gauge-coupled quantity, demonstrated by a four-way, Vicon-validated ablation) --- and even this is only \textbf{partially novel}. The underlying identifiability mathematics is already established in the acoustic / TDOA self-calibration literature \cite{kuang2013stratified,batstone2019robust} and in surveying error-propagation theory \cite{lichti2010propagation}; therefore only the empirical UWB demonstration and the operational thesis (``antenna-delay constants must not be transferred across coordinate frames'') are new. Contributions 2 (the integrated system), 3 (broadcast scheduling) and 5 (Monte-Carlo dropout / DOP-based redundancy) are established or standard practice, defensible only as validated engineering and characterization. Contribution 4 (burst-poll SS-TWR for motion error) is partially novel at best and carries a naming collision with the existing AltDS-TWR method \cite{neirynck2016alternative}. Recommendation: reposition the paper around Contribution 1's empirical and operational result, explicitly cede the identifiability theory to the acoustic/TDOA lineage, demote Contributions 2/3/5 to validated engineering, and rename Contribution 4. So positioned, the paper is a clean, citable measurement contribution publishable at \emph{Measurement} or \emph{IEEE TIM}.

\begin{table}[htbp]
\centering
\caption{Summary of verdicts}
\renewcommand{\arraystretch}{1.15}
\begin{tabular}{@{}p{0.31\linewidth}p{0.25\linewidth}p{0.34\linewidth}@{}}
\toprule
\textbf{Contribution} & \textbf{Verdict} & \textbf{Closest prior art} \\
\midrule
C1 --- Delay--scale gauge coupling (frame-dependent antenna delay) & PARTIALLY NOVEL & \cite{lichti2010propagation,kuang2013stratified,batstone2019robust} \\
C2 --- Integrated self-cal + joint delay system & ESTABLISHED-NOT-NOVEL & \cite{piavanini2022calibration,hamer2018selfcalibrating,lutz2019visual} \\
C3 --- Broadcast 8-anchor 10 Hz scheduling & STANDARD-PRACTICE & \cite{grosswindhager2019snaploc} \\
C4 --- Burst-poll SS-TWR (motion error) & PARTIALLY NOVEL (naming risk) & \cite{sidorenko2020error,dotlic2018ranging,neirynck2016alternative} \\
C5 --- Monte-Carlo keep-k / DOP redundancy & ESTABLISHED-NOT-NOVEL & \cite{albaidhani2021anchor,albaidhani2019anchor} \\
\bottomrule
\end{tabular}
\end{table}

\section{Contribution 1 --- Delay--layout coupling / frame-dependent antenna delay}

\verdict{PARTIALLY NOVEL.}

The claim is that, in range-only inter-anchor self-calibration, a constant additive antenna-delay term and the global geometry scale form a single coupled gauge degree of freedom. They are not separately identifiable from pairwise ranges alone, so an external metric constraint is required to fix the gauge; therefore antenna delay is frame-dependent, not a transferable hardware constant. The evidence is a four-way ablation: imposing externally-surveyed ground-truth anchor coordinates without re-estimating delays degrades tag localization (311 mm RMSE), while re-estimating delays in the imposed frame recovers and improves it (78 mm). The validation uses Vicon motion capture with SE(3)/Sim(3) alignment.

The closest prior art, ordered by disconfirming strength, is as follows. Lichti et al.\ \cite{lichti2010propagation} show that, in minimally constrained self-calibration, an unmodeled additive rangefinder offset produces near-linear scale-type patterns in range residuals, in a pulsed-radio/UWB context. This directly anticipates the additive-delay to scale coupling, although it does not frame it as a gauge freedom resolved by motion-capture ablation inside a tag-localization pipeline. Kuang and {\AA}str{\"o}m \cite{kuang2013stratified} and Kuang et al.\ \cite{kuang2013complete} are the most likely scooping lineage: TOA/TDOA-plus-unknown-offset is the mathematical isomorph of TWR-plus-antenna-delay, and those works characterize and resolve the affine/gauge ambiguities including the constant offset. Batstone et al.\ \cite{batstone2019robust} estimate sender/receiver positions under an unknown constant offset, which is the structural analogue of constant antenna delay coupled to geometry. Crocco et al.\ \cite{crocco2012closed} and Simayijiang et al.\ \cite{simayijiang2015toa} give closed-form microphone-array self-calibration methods that handle unknown internal delays/offsets, showing that joint geometry+offset self-calibration is mature. Shah et al.\ \cite{shah2021antenna} show that uncalibrated delay degrades localization even with correct positions, supporting the consequence, but they treat delay as a hardware constant, which is the assumption argued against. Shalaby et al.\ \cite{shalaby2023calibration} operationally treat antenna delay as calibrate-once/transferable, making it the key contrast reference for the operational thesis.

The contribution still survives because none of the acoustic/surveying works close the loop with a UWB TWR hardware realization, a four-way ablation isolating delay re-estimation in an imposed frame, and Vicon SE(3)/Sim(3) ground truth. The novelty should be positioned as empirical and operational, not as new identifiability mathematics.

\section{Contribution 2 --- Integrated system}

\verdict{ESTABLISHED-NOT-NOVEL as a concept; defensible only as a validated open implementation.}

The integrated system combines inter-anchor self-calibration, joint antenna-delay estimation, dual-layer anchor geometry, and custom broadcast SS-TWR firmware on DWM1001C/DW1000 while bypassing PANS/DRTLS. Piavanini et al.\ \cite{piavanini2022selfcalibrating,piavanini2022calibration} embed Gauss--Newton joint estimation of anchor geometry and antenna delays, which is the central pairing claimed here. Lutz et al.\ \cite{lutz2019visual} jointly estimate anchor poses and UWB sensor-error-model parameters. Hamer and D'Andrea \cite{hamer2018selfcalibrating} built a self-calibrating UWB network with clock synchronization, per-module self-localization, and simultaneous multi-robot operation via TDOA, with no central coordination. Nguyen et al.\ \cite{nguyen2025coordinate} jointly estimate 3D anchor positions and range biases via continuous-time factor-graph calibration. Custom DW1000 broadcast firmware bypassing PANS/DRTLS is widespread.

The defensible angle is that the specific assembly --- dual-layer geometry, broadcast SS-TWR custom firmware on DWM1001C, a gauge-aware workflow, and released Vicon-validated evidence --- is a credible but incremental engineering contribution. It should be claimed as a validated, reproducible system, not as conceptual novelty.

\section{Contribution 3 --- Broadcast 8-anchor 10 Hz scheduling}

\verdict{STANDARD-PRACTICE. Confirmed.}

The claimed scheduling result is 8999 fixes/300 s with a balanced 3-tag load, zero timeouts, and a runtime-reconfigurable layout. Broadcast/scalable operation is inherent to TDoA RTLS and to SnapLoc \cite{grosswindhager2019snaploc}, which lets an unlimited number of tags self-localize at a theoretical upper bound of 2.3 kHz with decimeter accuracy on the off-the-shelf DW1000. One-to-many ranging is codified in IEEE 802.15.4z. The throughput/reliability numbers are a useful reliability and throughput characterization, not conceptual novelty.

\section{Contribution 4 --- ``Alt SS-TWR burst-poll'' for motion-error reduction (\texorpdfstring{\textasciitilde{}8$\times$}{~8x})}

\verdict{PARTIALLY NOVEL; high naming-collision risk; needs sharper framing.}

Sidorenko et al.\ \cite{sidorenko2020error} are the key reference to beat: they model constant-velocity motion as a linear pseudo-clock-drift and correct it, pre-empting the conceptual mechanism. Dotlic et al.\ \cite{dotlic2018ranging} give CFO-corrected SS-TWR with fewer frames, overlapping the fast-SS-TWR-with-reduced-error goal. Shalaby et al.\ \cite{shalaby2024reducing} reduce variance via signal-timing optimization, overlapping the timing-design angle. Neirynck et al.\ \cite{neirynck2016alternative} create the naming collision: ``AltDS-TWR'' already denotes a specific double-sided method whose error depends only on one device's clock drift, and the term is now standard. The proposed single-sided burst-poll must be renamed in the manuscript.

A burst-poll SS-TWR specifically targeting motion-induced error with a quantified \textasciitilde{}8$\times$ reduction is not exactly anticipated, but the mechanism is well-trodden. Novelty is marginal without a clean theoretical model distinguishing it from Sidorenko et al.\ \cite{sidorenko2020error}.

\section{Contribution 5 --- Monte-Carlo keep-k dropout + stratified vertical-coverage + anchor-redundancy-to-accuracy guarantee}

\verdict{ESTABLISHED-NOT-NOVEL. Confirmed (DOP/CRLB territory).}

\begin{sloppypar}
Anchor-redundancy-to-accuracy and geometry-quality analysis are classical DOP/GDOP/WGDOP\allowbreak{} and CRLB results. Albaidhani and Alsudani \cite{albaidhani2021anchor} address anchor selection by GDOP for UWB, and Albaidhani et al.\ \cite{albaidhani2019anchor} address anchor selection via MSE/WGDOP. k-coverage is standard in sensor networks; the per-fix accuracy-guarantee framing mirrors GNSS integrity monitoring, RAIM, and protection levels; and the stratified vertical-coverage finding is the well-known z-axis GDOP effect. The Monte-Carlo keep-k study is a reasonable empirical robustness characterization, not a new method.
\end{sloppypar}

\section{Conclusions}

\paragraph{(a) Defensible single-sentence novelty statement.}
\begin{displayquote}
``We provide the first explicit empirical demonstration, validated against motion-capture ground truth, that in range-only UWB inter-anchor self-calibration the additive antenna delay and the global geometric scale constitute a single non-identifiable gauge degree of freedom --- so antenna delay must be re-estimated within whatever metric frame is imposed and is therefore not a transferable hardware constant.''
\end{displayquote}

\paragraph{(b) Experiments that would strengthen each thin claim.}
\begin{itemize}[leftmargin=*]
\item C1: add an explicit observability / Fisher-information (rank-deficiency) analysis of the delay--scale Jacobian, exhibit the degenerate singular value / null-space direction numerically, cite the acoustic isomorphism \cite{kuang2013stratified,batstone2019robust} and \cite{lichti2010propagation} directly, and replicate the 311\(\to\)78 mm ablation across \(\geq\)2 independent anchor geometries.
\item C2: head-to-head quantitative comparison against \cite{piavanini2022calibration} and \cite{hamer2018selfcalibrating} on the same dataset; release firmware + dataset.
\item C3: reframe as a reliability/latency benchmark vs \cite{grosswindhager2019snaploc} / a PANS-TDoA baseline.
\item C4: derive a closed-form motion-error model for burst-poll SS-TWR and compare to \cite{sidorenko2020error} and \cite{dotlic2018ranging} on identical motion profiles; rename the scheme.
\item C5: replace/augment the Monte-Carlo study with a CRLB or GDOP bound and show the keep-k accuracy curve tracks the bound.
\end{itemize}

\paragraph{(c) Overall positioning.}
The work is publishable at \emph{Measurement} or \emph{IEEE TIM} if and only if repositioned around Contribution 1's empirical and operational result, explicitly ceding the identifiability mathematics to the acoustic/TDOA self-calibration and surveying literature, demoting Contributions 2--3 and 5 to validated engineering/characterization, and strengthening or renaming Contribution 4. Claiming broad theoretical novelty for the gauge coupling would invite a strong reject from any reviewer familiar with the \cite{kuang2013stratified} / \cite{batstone2019robust} lineage. The operational thesis (``do not transfer antenna-delay constants across coordinate frames / deployments''), backed by the ablation, is a clean, citable, useful result --- that is where the paper's value lives.

\begin{thebibliography}{99}

\bibitem{lichti2010propagation}
D. D. Lichti, K. O'Keefe, and S. Jamtsho, ``Propagation of an Unmodeled Additive Constant in Range Sensor Observations,'' \emph{Journal of Surveying Engineering}, 2010, vol. 136, no. 3, pp. 111--119, doi: 10.1061/(ASCE)SU.1943-5428.0000024.

\bibitem{kuang2013stratified}
Yubin Kuang and Kalle {\AA}str{\"o}m, ``Stratified Sensor Network Self-Calibration from TDOA Measurements,'' in \emph{21st European Signal Processing Conference (EUSIPCO 2013)}, 2013, pp. 1--5, IEEE, note: verify page range.

\bibitem{kuang2013complete}
Yubin Kuang, Simon Burgess, Anna Torstensson, and Kalle {\AA}str{\"o}m, ``A Complete Characterization and Solution to the Microphone Position Self-Calibration Problem,'' in \emph{2013 IEEE International Conference on Acoustics, Speech and Signal Processing (ICASSP)}, 2013, pp. 3875--3879, IEEE, doi: 10.1109/ICASSP.2013.6638382.

\bibitem{batstone2019robust}
Kenneth Batstone, Gabrielle Flood, Thejasvi Beleyur, Viktor Larsson, Holger R. Goerlitz, Magnus Oskarsson, and Kalle {\AA}str{\"o}m, ``Robust Self-Calibration of Constant Offset Time-Difference-of-Arrival,'' in \emph{2019 IEEE International Conference on Acoustics, Speech and Signal Processing (ICASSP)}, 2019, pp. 4410--4414, IEEE, doi: 10.1109/ICASSP.2019.8683221.

\bibitem{crocco2012closed}
Marco Crocco, Alessio Del Bue, Marco Bustreo, and Vittorio Murino, ``A Closed Form Solution to the Microphone Position Self-Calibration Problem,'' in \emph{2012 IEEE International Conference on Acoustics, Speech and Signal Processing (ICASSP)}, 2012, pp. 2597--2600, IEEE, doi: 10.1109/ICASSP.2012.6288448.

\bibitem{simayijiang2015toa}
Zhayida Simayijiang, Simon Burgess, Yubin Kuang, and Kalle {\AA}str{\"o}m, ``TOA-Based Self-Calibration of Dual-Microphone Array,'' \emph{IEEE Journal of Selected Topics in Signal Processing}, 2015, vol. 9, no. 5, pp. 791--801, doi: 10.1109/JSTSP.2015.2417117.

\bibitem{shah2021antenna}
Syed Aamir Hussain Shah and others, ``Antenna Delay Calibration of UWB Nodes,'' \emph{IEEE Access}, 2021, vol. 9, pp. 63294--63305, doi: 10.1109/ACCESS.2021.3074872, note: verify full author list, volume and pages.

\bibitem{shalaby2023calibration}
Mohammed Ayman Shalaby, Charles Champagne Cossette, James Richard Forbes, and Jerome Le Ny, ``Calibration and Uncertainty Characterization for Ultra-Wideband Two-Way-Ranging Measurements,'' in \emph{2023 IEEE International Conference on Robotics and Automation (ICRA)}, 2023, pp. 4128--4134, IEEE, doi: 10.1109/ICRA48891.2023.10160769.

\bibitem{piavanini2022selfcalibrating}
Marco Piavanini, Luca Barbieri, Mattia Brambilla, Mattia Cerutti, Simone Ercoli, Andrea Agili, and Monica Nicoli, ``A Self-Calibrating Localization Solution for Sport Applications with UWB Technology,'' \emph{Sensors}, 2022, vol. 22, no. 23, p. 9363, doi: 10.3390/s22239363.

\bibitem{piavanini2022calibration}
Marco Piavanini, Luca Barbieri, Mattia Brambilla, Mattia Cerutti, Simone Ercoli, Andrea Agili, and Monica Nicoli, ``A Calibration Method for Antenna Delay Estimation and Anchor Self-Localization in UWB Systems,'' in \emph{2022 IEEE International Workshop on Metrology for Industry 4.0 \& IoT (MetroInd4.0\&IoT)}, Trento, Italy, 2022, pp. 173--177, IEEE, doi: 10.1109/MetroInd4.0IoT54413.2022.9831579.

\bibitem{hamer2018selfcalibrating}
Michael Hamer and Raffaello D'Andrea, ``Self-Calibrating Ultra-Wideband Network Supporting Multi-Robot Localization,'' \emph{IEEE Access}, 2018, vol. 6, pp. 22292--22304, doi: 10.1109/ACCESS.2018.2829020.

\bibitem{nguyen2025coordinate}
Tien-Dat Nguyen, Thien-Minh Nguyen, and Vinh-Hao Nguyen, ``Coordinate-Consistent Localization via Continuous-Time Calibration and Fusion of UWB and SLAM Observations,'' 2025, arXiv:2510.05992, cs.RO.

\bibitem{grosswindhager2019snaploc}
Bernhard Gro{\ss}windhager, Michael Stocker, Michael Rath, Carlo Alberto Boano, and Kay R{\"o}mer, ``SnapLoc: An Ultra-Fast UWB-Based Indoor Localization System for an Unlimited Number of Tags,'' in \emph{2019 18th ACM/IEEE International Conference on Information Processing in Sensor Networks (IPSN)}, 2019, pp. 61--72, IEEE, doi: 10.1145/3302506.3310389.

\bibitem{sidorenko2020error}
Juri Sidorenko, Volker Schatz, Norbert Scherer-Negenborn, Michael Arens, and Urs Hugentobler, ``Error Corrections for Ultrawideband Ranging,'' \emph{IEEE Transactions on Instrumentation and Measurement}, 2020, vol. 69, no. 11, pp. 9037--9047, doi: 10.1109/TIM.2020.2996706.

\bibitem{dotlic2018ranging}
Igor Dotlic, Andrew Connell, and Michael McLaughlin, ``Ranging Methods Utilizing Carrier Frequency Offset Estimation,'' in \emph{2018 15th Workshop on Positioning, Navigation and Communications (WPNC)}, 2018, pp. 1--6, IEEE, doi: 10.1109/WPNC.2018.8555763.

\bibitem{neirynck2016alternative}
Dries Neirynck, Eric Luk, and Michael McLaughlin, ``An Alternative Double-Sided Two-Way Ranging Method,'' in \emph{2016 13th Workshop on Positioning, Navigation and Communications (WPNC)}, 2016, pp. 1--4, IEEE, doi: 10.1109/WPNC.2016.7822844.

\bibitem{shalaby2024reducing}
Mohammed Ayman Shalaby, Charles Champagne Cossette, James Richard Forbes, and Jerome Le Ny, ``Reducing Two-Way Ranging Variance by Signal-Timing Optimization,'' \emph{IEEE Transactions on Aerospace and Electronic Systems}, 2024, arXiv:2211.00538, note: verify volume/number/pages on IEEE Xplore.

\bibitem{albaidhani2021anchor}
Abbas Albaidhani and Ahmed Alsudani, ``Anchor Selection by Geometric Dilution of Precision for an Indoor Positioning System Using Ultra-Wide Band Technology,'' \emph{IET Wireless Sensor Systems}, 2021, vol. 11, no. 1, pp. 22--31, doi: 10.1049/wss2.12006.

\bibitem{albaidhani2019anchor}
Abbas Albaidhani, Antoni Morell, and Jos{\'e} L{\'o}pez Vicario, ``Anchor Selection for UWB Indoor Positioning,'' \emph{Transactions on Emerging Telecommunications Technologies}, 2019, vol. 30, no. 6, e3598, doi: 10.1002/ett.3598.

\bibitem{lutz2019visual}
Philipp Lutz, Martin J. Schuster, and Florian Steidle, ``Visual-Inertial SLAM Aided Estimation of Anchor Poses and Sensor Error Model Parameters of UWB Radio Modules,'' in \emph{2019 19th International Conference on Advanced Robotics (ICAR)}, 2019, pp. 739--746, IEEE, doi: 10.1109/ICAR46387.2019.8981544, note: verify page range.

\end{thebibliography}

\end{document}
